\section{Discussion}

\subsection{Model Performance Patterns}

Our evaluation was designed to test whether frontier LLMs reliably exceed the USMLE passing threshold and whether meaningful performance differences exist among commercially available models. Based on prior studies of GPT-4 \parencite{nori2023gpt4usmle} and Claude \parencite{anthropic2024claude3}, we anticipated that the largest models from OpenAI, Anthropic, and Google would demonstrate the highest accuracy, with smaller variants (GPT-4o-mini, Claude 3 Haiku, Gemini 1.5 Flash) showing modest reductions. We also hypothesized that DeepSeek R1 Distill and Llama 3.3 70B, despite being open-weight models running via a third-party inference provider (Groq), would compete favorably on Step 1 basic science questions—where the reasoning chains are more formulaic—but show larger drops on Step 2 CK clinical vignettes, which require integrating patient presentation details, recognizing rare disease patterns, and navigating nuanced management tradeoffs.

If the pattern of Step 2 CK being more challenging than Step 1 is confirmed empirically, the implication is significant: students relying on LLMs for Step 2 CK preparation—the more clinically consequential examination for direct patient care—may receive less reliable assistance than they would for Step 1 review. This is particularly concerning given that LLM adoption appears to be especially high among students in the clinical years, where Step 2 CK preparation is most active.

\subsection{Subject-Level Weaknesses}

Subject-level analysis is likely to reveal heterogeneous model performance. We anticipate relatively strong performance in domains with abundant training data and clearly defined right-and-wrong answers (e.g., Biochemistry, Pharmacology, Microbiology), and relatively weaker performance in domains requiring contextual clinical judgment (e.g., Psychiatry, Behavioral Science, Biostatistics). Such patterns would suggest that LLMs are safer to rely on for foundational science content than for complex clinical or psychosocial reasoning—an important nuance for study strategy.

\subsection{Implications for IMG Students}

The IMG perspective analysis is the most novel contribution of this work. If our hypothesis holds—that models show lower accuracy on questions encoding US-centric healthcare system assumptions—the practical implication is that IMG students using these tools for Step 2 CK preparation may receive systematically lower-quality guidance on precisely the question types that most frequently disadvantage IMG candidates. This would represent a compounding disadvantage: the examination is already harder for IMGs due to unfamiliar clinical contexts, and the AI study tool may be less reliable on those same questions.

Even if the accuracy gap is not statistically significant in our relatively small sample (200 questions), any observed directional trend warrants caution and motivates a larger confirmatory study. IMG students should be informed that LLM accuracy on US-centric questions has not been validated and that supplementing AI tools with US-specific clinical resources (e.g., UWorld, Amboss, First Aid for the USMLE Step 2 CK) remains important.

\subsection{Limitations}

Several limitations should be acknowledged. First, our sample of 200 questions is modest for detecting small accuracy differences with high statistical power; a larger confirmatory study using the full MedQA dataset or retired USMLE items would strengthen the conclusions. Second, MedQA, while widely used as a research benchmark, may not fully represent the current difficulty, distribution, or item formats of actual USMLE examinations, which are updated regularly by the NBME. Third, we employed zero-shot prompting with a fixed chain-of-thought template; actual student use of LLMs is more varied and often involves follow-up questions, multi-turn dialogue, and retrieval of supplementary material. Our results reflect one specific prompting regime. Fourth, the keyword-based IMG-relevance tagger, while validated by our medical co-author, is a rough proxy; a gold-standard annotation of all 200 questions would provide a more precise estimate of the IMG accuracy gap. Fifth, the expert annotation sample of 50 responses per model provides useful qualitative context but limited statistical power for detecting differences in reasoning quality.

\subsection{Ethical Considerations}

The use of AI tools as study aids raises ethical questions about equity, over-reliance, and clinical safety. Students who cannot afford commercial question banks may over-rely on LLMs, accepting confident but incorrect answers without the calibration that comes from knowing a tool's track record on a validated item pool. Models that hallucinate with apparent authority pose a particular risk when the user lacks the domain knowledge to recognize errors. We recommend that LLMs be used as a \textit{supplement} to—not a replacement for—structured studying with validated, NBME-aligned question banks, and that any AI-generated clinical reasoning be critically evaluated rather than accepted uncritically.

\subsection{Future Directions}

This study opens several directions for future research. First, evaluating multi-turn prompting strategies (e.g., asking the model to reconsider its answer or to explain a specific step in its reasoning) may reveal accuracy improvements and is more representative of actual student use. Second, retrieval-augmented generation (RAG) approaches that ground model responses in authoritative medical references (e.g., UpToDate, Harrison's Principles) may substantially reduce hallucinations on clinical vignettes. Third, extending the evaluation to non-English prompting—specifically Hindi, Spanish, and Mandarin, which are among the most common first languages of IMG students—would directly assess whether language of inquiry affects accuracy on the English-language examination. Finally, collaboration with the NBME to evaluate models on retired, psychometrically validated USMLE items would provide a more definitive benchmark than is currently possible with publicly available question banks.
